\section{Attainable loads and a regularized inverse}\label{sec:reachability}

\subsection{Continuous quadratic load operators}
At fixed geometry let $\mathcal A_\Gamma$ denote a well-posed linear acoustic
operator including its boundary and interface conditions, and let
$\mathcal B_\Gamma$ map a source field to its right-hand side. The acoustic
state $\bm z_a$ satisfies
\begin{equation}
 \mathcal A_\Gamma\bm z_a=\mathcal B_\Gamma u.
 \label{eq:distributed-state}
\end{equation}
Let $\mathcal S_\Gamma:\mathcal U\to\mathcal X$ map sources to the pressure
and velocity traces needed for the radiation traction. Here $\mathcal X$ is
a Hilbert trace space, with fixed physical scaling, continuously embedded in
the product of $L^4(\Gamma)$ spaces for those components. The latter space
contains functions with integrable fourth power. This regularity ensures that
quadratic loads lie in $L^2(\Gamma)$. Such trace regularity is an assumption
on this reduction, not a property of every sharp-edged or singular-source
configuration.

For any real dimensionless test function $\phi$ compatible with the surface
support and with $\int_\Gamma\phi\dd A=0$, define the generalized force
and the required force by
\begin{equation}
 F_\phi(u)=\int_\Gamma\phi\Pi(u)\dd A,
 \qquad b_\phi=\int_\Gamma\phi(\sigma\kappa+\Delta\rho\Phi)\dd A.
 \label{eq:continuous-force}
\end{equation}
Both have units of newtons. Bounded tests suffice; more general tests are
allowed where the pairing is continuous. The radiation expression defines a
bounded Hermitian quadratic form $\mathcal C_\phi$ on $\mathcal X$. With
the star denoting the Hilbert adjoint in this subsection, set
\begin{equation}
 \mathcal K_\phi=\mathcal S_\Gamma^*\mathcal C_\phi\mathcal S_\Gamma,
 \qquad
 \boxed{F_\phi(u)=\inner{u}{\mathcal K_\phi u}_{\mathcal U}=b_\phi.}
 \label{eq:continuous-quadratic}
\end{equation}
The operator $\mathcal K_\phi$ is self-adjoint, although it need not be
positive. Because the source norm is dimensionless, its quadratic form has
force units. Exact quiescent normal balance requires this equality for a
dense set of admissible tests, with the wetting and volume constraints.
Tangential balance and absence of incompatible mean bulk forcing remain
separate requirements. A steady streaming state instead uses the full
stationary equations.

At fixed geometry the first variation is
\begin{equation}
 DF_\phi(u)[\delta u]=
       2\Rea\inner{\mathcal K_\phi u}{\delta u}_{\mathcal U}.
 \label{eq:continuous-load-derivative}
\end{equation}
Multiplication of all sources by the same phase $e^{\mathrm i\vartheta}$,
for real $\vartheta$, leaves the load unchanged. Consequently $\delta u=
\mathrm i u\,\delta\vartheta$ is a null direction of this derivative.
Sources invisible to the relevant wave traces add further nonuniqueness.
At $u=0$ the derivative vanishes; this does not exclude finite-amplitude
shaping. For infinitesimal target changes, the differentiated required load,
including the change of $\mathcal K_\phi$ with geometry, must lie in the
admissible derivative range. This is a condition of local authority, not a
prescription to generate the family by continuation.

\subsection{Necessary source-effort and sign bounds}
Choose a finite set of tests $\phi_i$, real coefficients $a_i$, and write
$b_i=b_{\phi_i}$, $\mathcal K_i=\mathcal K_{\phi_i}$ and
$\mathcal K_a=\sum_i a_i\mathcal K_i$. Let $C_u>0$ be an upper bound on
$\norm{u}_{\mathcal U}$ and let $\norm{\mathcal K_a}$ denote operator norm.
Every exact source under that bound satisfies
\begin{equation}
 \left|\sum_i a_i b_i\right|\leq C_u^2\norm{\mathcal K_a},
 \qquad
 \mathcal K_a\succeq0\ \Longrightarrow\ \sum_i a_i b_i\geq0.
 \label{eq:continuous-certificates}
\end{equation}
Here $\succeq0$ means a nonnegative quadratic form. The inequalities follow
directly from \cref{eq:continuous-quadratic} and the operator-norm bound.
If $\norm{\mathcal K_a}>0$, they also give the lower bound
$\norm{u}_{\mathcal U}^2\geq|\sum_i a_i b_i|/\norm{\mathcal K_a}$.
If the operator is zero but the target combination is nonzero, exact matching
is impossible in this model. These are certificates of necessary conditions;
passing them does not establish reachability, stability or uniqueness.

A power version needs an additional hypothesis. Suppose input power is the
quadratic form $\inner{u}{\mathcal R_Pu}_{\mathcal U}$ with a bounded,
coercive self-adjoint operator $\mathcal R_P$, meaning
$\inner{u}{\mathcal R_Pu}\geq r_P\norm{u}_{\mathcal U}^2$ for some
$r_P>0$. If $P_{\max}$ bounds that power, then
\begin{equation}
 \left|\sum_i a_i b_i\right|
 \leq P_{\max}\,
 \norm{\mathcal R_P^{-1/2}\mathcal K_a\mathcal R_P^{-1/2}}.
 \label{eq:continuous-power-bound}
\end{equation}
This follows by substituting $\mathcal R_P^{1/2}u$ in the quadratic form.
Passivity alone does not give coercivity: a nearly lossless resonant source
can have very different stored energy, amplitude and net consumed power.
The source-effort bound is therefore not silently relabeled a power bound.

\subsection{Why arbitrary source functions do not imply arbitrary loads}
Let $\mathcal Y=L^2_0(\Gamma)$ be the real space of square-integrable
zero-mean loads, equipped with the dimensionless norm
$\norm{f}_{\mathcal Y}^2=\int_\Gamma |f|^2\dd A/(A[\Gamma]\Pi_0^2)$,
where $\Pi_0>0$ is a chosen stress scale. Define the continuous quadratic
map $\mathcal Q:\mathcal X\to\mathcal Y$ by taking radiation traction
and applying the zero-mean projection $\mathcal P_0$ of
\cref{eq:pressure-gauge}. The required load in this space is
$b=\mathcal P_0(\sigma\kappa+\Delta\rho\Phi)$.

\begin{proposition}[A conditional obstruction under bounded source effort]
Suppose $\mathcal U$ is Hilbert, $\mathcal S_\Gamma:\mathcal U\to\mathcal X$
is compact and linear, and $\mathcal Q$ is continuous. The set
\begin{equation}
 \mathfrak R_{C_u}=\{\mathcal Q(\mathcal S_\Gamma u):
                         \norm{u}_{\mathcal U}\leq C_u\}
 \label{eq:bounded-reachable-set}
\end{equation}
is compact in $\mathcal Y$. If $\mathcal Y$ is infinite dimensional, this
set has empty interior in $\mathcal Y$.
\end{proposition}
\begin{proof}
Any sequence in the closed source ball has a weakly convergent subsequence
with limit in that ball. Compactness of $\mathcal S_\Gamma$ makes the
corresponding traces converge strongly. Continuity of $\mathcal Q$ then
makes the loads converge to the load of the limiting source. Thus the image
is sequentially compact, hence compact in the normed load space. A compact
subset of an infinite-dimensional normed space cannot contain an open ball:
its closure would contain a scaled closed unit ball, which is not compact.
\end{proof}

The hypothesis expresses spatial smoothing between regular sources and the
interface. Positive source-interface separation, regular geometry and a
frequency away from unresolved poles are settings in which such smoothing
can be investigated; cylinder geometry alone does not prove compactness.
Sources acting directly on the observed interface, singular distributions,
changing source supports or an unbounded frequency range are not automatically
covered. Nor does the result rule out a finite-dimensional Cartesian target
family. It rules out claiming that bounded distributed actuation covers an
open set of unrestricted load fields merely because the source has infinitely
many degrees of freedom. This is a deduction for the stated operator, not a
published universal no-go theorem about acoustic lenses.

\subsection{A well-defined regularized fixed-target problem}
Let $\gamma_u>0$ be a dimensionless regularization weight. A reusable
fixed-geometry inverse is
\begin{equation}
 \min_{u\in\mathcal U_{\rm ad}} j_\Gamma(u),\qquad
 j_\Gamma(u)=\frac12\norm{\mathcal Q(\mathcal S_\Gamma u)-b}_{\mathcal Y}^2
                 +\frac{\gamma_u}{2}\norm{u}_{\mathcal U}^2.
 \label{eq:regularized-source-inverse}
\end{equation}
It minimizes load error and declared source effort. The positive penalty
selects controlled effort; it may deliberately leave nonzero mismatch even
when an exact match exists. A tolerance-constrained minimum-effort problem
is another choice. Neither is generally convex because the load is quadratic
in coherent source amplitude.

\begin{proposition}[Conditional existence of a regularized source]
Under the compactness assumptions above, if $\mathcal U_{\rm ad}$ is
nonempty and weakly closed, \cref{eq:regularized-source-inverse} has a
minimizer.
\end{proposition}
\begin{proof}
A minimizing sequence has bounded norm because $\gamma_u>0$ and the
mismatch is nonnegative. Select a weakly convergent subsequence in
$\mathcal U$. Its limit is admissible by weak closedness. Compact wave
transfer and continuous quadratic loading give strong convergence of the
load, while the squared source norm is weakly lower semicontinuous. The
limit therefore attains the infimum.
\end{proof}
This proof guarantees a minimizer of the declared reduced problem, not zero
error, uniqueness, an algorithm finding the global minimum, a stable fluid
state, or existence for the full moving-interface system. Nonconvex hardware
constraints need their own weak-closedness check. Regularization and support
restrictions are also central to linear Helmholtz source problems
\cite{hubenthal2016,pieper2020}; the nonlinear load problem above has its own
explicit assumptions and proof.

\subsection{A second-order entry that needs no neighboring target}
A concrete analytic construction addresses the zero-source degeneracy. Assume
zero actuation is admissible. Choose
$M$ generalized forces and positive scales $F_{i,0}$ in newtons. For a
positive dimensionless penalty $\gamma_u$, consider the projected cost
\begin{equation}
 j_M(u)=\frac12\sum_{i=1}^M
       \left(\frac{F_i(u)-b_i}{F_{i,0}}\right)^2
                      +\frac{\gamma_u}{2}\norm{u}_{\mathcal U}^2.
 \label{eq:projected-source-cost}
\end{equation}
Define the dimensionless self-adjoint operator
\begin{equation}
 \mathcal H_0=\gamma_u\mathcal I_{\mathcal U}
                   -2\sum_i\frac{b_i}{F_{i,0}^2}\mathcal K_i,
 \label{eq:zero-source-hessian}
\end{equation}
where $\mathcal I_{\mathcal U}$ is the identity on the source space. Expanding
the squares gives the exact identity
\begin{equation}
 j_M(u)-j_M(0)=\frac12\inner{u}{\mathcal H_0u}_{\mathcal U}
                 +\frac12\sum_i\left(\frac{F_i(u)}{F_{i,0}}\right)^2.
 \label{eq:zero-source-expansion}
\end{equation}
Thus $\mathcal H_0\succeq0$ makes zero a global minimizer of this particular
penalized objective, even though it need not meet the target tolerance. This
signals a limitation of the objective or available wave response, not proof
that the target is physically impossible. If an admissible direction has
negative quadratic form, a sufficiently small nonzero source in that direction
strictly improves the cost. The real first derivative at zero cannot detect it.

There is also an exact amplitude optimization along any direction. Let
$d\in\mathcal U$ have unit norm, write $u=\sqrt{x}\,d$ with $x\geq0$,
and define $f_i=F_i(d)/F_{i,0}$, $\widehat b_i=b_i/F_{i,0}$,
$A_d=\sum_i\widehat b_i f_i-\gamma_u/2$, and $B_d=\sum_i f_i^2$.
Then
\begin{equation}
 j_M(\sqrt{x}\,d)=j_M(0)-A_dx+\frac12B_dx^2,
 \qquad
 x_\star=\min\{C_u^2,\max\{0,A_d/B_d\}\}
 \quad (B_d>0),
 \label{eq:source-ray-amplitude}
\end{equation}
for the sole bound $\norm{u}_{\mathcal U}\leq C_u$. If $B_d=0$, the positive
regularization makes $x_\star=0$. Additional physical constraints change the
admissible interval or can invalidate the direction. A negative-curvature
direction of $\mathcal H_0$ and this amplitude rule give a principled entry
to a source solve for the target at hand. Optimizing the direction, enforcing
the full unsampled load, and establishing fluid stability remain necessary;
this construction is not a global inverse solution or a family-coverage proof.

\subsection{What changes when the surface is also unknown}
Let $\bm y$ collect the surface, hydrodynamic state, acoustic state and, where
needed, temperature and actuator variables. Let
$\mathcal R(\bm y,u)=0$ collect the stationary governing equations,
volume conditions, boundary conditions and chosen pressure gauges.
Let $\Gamma_{\rm tar}\subseteq\Gamma_\star$ denote the target patch, of area
$A_{\rm tar}>0$. For a nearby surface let $\eta_s$ be its signed normal
departure on this fixed target patch, and let $h_0>0$ be a geometric tolerance.
Define the dimensionless shape error
\begin{equation}
 J_s=\frac{1}{A_{\rm tar}h_0^2}\int_{\Gamma_{\rm tar}}\eta_s^2\dd A.
 \label{eq:shape-cost}
\end{equation}
Larger deformations need a declared surface-distance and correspondence
convention. Let $J_{\rm obs}\geq0$ denote any additional dimensionless
application-specific observation error, and choose nonnegative weights
$\beta_s,\beta_o,\beta_u$. A coupled design can minimize
\begin{equation}
 J_{\rm stat}=\beta_sJ_s+\beta_oJ_{\rm obs}
                         +\beta_u\norm{u}_{\mathcal U}^2
 \quad\text{subject to }\mathcal R(\bm y,u)=0.
 \label{eq:coupled-static-design}
\end{equation}
The admissible set also imposes the declared source and stability limits. For
an exact target impose $\Gamma=\Gamma_\star$ directly; a soft penalty alone
does not enforce exact shape. Optical observations are a later specialization.
The full problem is nonlinear
because both required curvature and acoustic transfer depend on geometry;
the fixed-target relaxation does not make it globally convex.

Let $\delta$ denote a first variation of geometry, material state or drive.
Differentiating \cref{eq:distributed-state} gives
\begin{equation}
 \mathcal A_\Gamma\,\delta\bm z_a
   =(\delta\mathcal B_\Gamma)u
       +\mathcal B_\Gamma\delta u
       -(\delta\mathcal A_\Gamma)\bm z_a.
 \label{eq:wave-sensitivity}
\end{equation}
Moving normals, surface measure, interface conditions and material coefficients
are part of this derivative. Freezing the acoustic field while differentiating
only capillary curvature omits the central wave--shape feedback.

If the linearized constrained stationary operator $\mathcal R_{\bm y}$ is
invertible on the gauge-fixed state space, implicit differentiation gives
\begin{equation}
 \delta\bm y=-\mathcal R_{\bm y}^{-1}
                         \mathcal R_u\,\delta u.
 \label{eq:implicit-sensitivity}
\end{equation}
At a fold or singular branch that inverse fails, and a local sensitivity cannot
be interpreted as a global design map. Resonance can also make it very poorly
conditioned.

Finally, let $\bm a$ be the actually commanded electrical or mechanical
field and let $u=\mathcal H(\Gamma,\omega)\bm a$ be its loaded source response.
Holding $\bm a$ fixed generally does not hold $u$ fixed. The shape derivative
includes $\delta\mathcal H$. Prescribed boundary velocity is experimentally
relevant only when it is controlled or when the transducer dynamics needed to
realize it are included. The same distinction applies to source gradients.

\subsection{Finite arrays approximate a source space}
Let $\mathcal U_N$ be a nested sequence of finite-dimensional source spaces.
An array or basis represents $u=E_N\bm w$, where $E_N$ maps its $N$ complex
coefficients to physical source fields. This is a restriction of
\cref{eq:regularized-source-inverse}, not its definition. Assume that every
$u\in\mathcal U_{\rm ad}$ admits strongly convergent admissible
approximants $u_N\in\mathcal U_N\cap\mathcal U_{\rm ad}$.

\begin{proposition}[Conditional approximation of the inverse value]
Under the existence hypotheses and the admissible-approximation assumption,
the minima of $j_\Gamma$ over $\mathcal U_N\cap\mathcal U_{\rm ad}$
converge to the minimum over $\mathcal U_{\rm ad}$. Every weak accumulation
point of corresponding minimizers is a minimizer of the continuous problem.
\end{proposition}
\begin{proof}
The restricted infima cannot be smaller than the continuous infimum. Strong
admissible approximants to a continuous minimizer have convergent objective
values, giving the reverse limiting bound. Their bounded values bound the
norms of the restricted minimizers. Apply the weak-compactness and lower
semicontinuity argument of the preceding proof to their subsequences.
\end{proof}
Density of an unconstrained basis alone is insufficient; the approximants
must respect the physical bounds. Arbitrary unrelated arrays need not form
such a sequence. This statement also assumes the exact same wave model; it
does not replace a separate convergence analysis of acoustic or fluid meshes.
